\documentclass[12pt,a4paper,parskip=half]{scrartcl}

%   Encoding & Fonts
\usepackage[utf8]{inputenc}
\usepackage[T1]{fontenc}
\usepackage{lmodern}
\usepackage{microtype}

%   Math
\usepackage{amsmath, amssymb, amsthm, mathtools}
\usepackage{bm}
\usepackage{physics}

%   Graphics
\usepackage{tikz}
\usepackage{pgfplots}
\pgfplotsset{compat=1.18}

%   Colors
\usepackage{xcolor}
\definecolor{alertred}{RGB}{180,0,0}

%   References
\usepackage[backend=biber, style=authoryear, sorting=nty]{biblatex}
\addbibresource{reference.bib}
\usepackage{hyperref}
\usepackage{cleveref}

%   Layout
\usepackage{geometry}
\geometry{a4paper, left=2.5cm, right=2.5cm, top=2.5cm, bottom=2.5cm}
\usepackage{setspace}
\onehalfspacing

%   Simple block environments (Standard-LaTeX, fußnotenfreundlich)
\newenvironment{block}[1]{%
  \begin{quote}
  \textbf{#1}\\[0.5em]
  \ignorespaces
}{%
  \end{quote}
}

\newenvironment{alertblock}[1]{%
  \begin{quote}
  \textcolor{alertred}{\textbf{#1}}\\[0.5em]
  \ignorespaces
}{%
  \end{quote}
}

%   Beamer-ähnlicher \alert-Befehl
\newcommand{\alert}[1]{\textcolor{alertred}{\textbf{#1}}}

%   Document Metadata
\title{Binary Payment Schemes: Moral Hazard and Loss Aversion\\[0.5em]
\large A Literature Review}
\author{Theo Osterhaus \\ {\normalsize Matriculation Number: 7443693}}
\date{Seminar Market Design and Behaviour, SS 2026\\[0.2em]
Department of Economics\\ Universität zu Köln}

\begin{document}

\maketitle

\tableofcontents
\clearpage

\section{Motivation}

\subsection*{Contractual Form Matters: A Problem of Decades}

\begin{itemize}
    \item \textbf{Observation:} Real-world contracts are often remarkably \textbf{simple} (e.g., binary bonus schemes, flat wages).
    \item \textbf{Classical Contract Theory} \parencite{Holmstroem1979}: Optimal contracts should be \textbf{complex}, exploiting all available information.
\end{itemize}

\begin{block}{The Central Puzzle}
    If contractual form matters so much, why do we observe such a prevalence of relatively simple contracts?
\end{block}

\begin{itemize}
    \item \textcite{Prendergast1999} documented the discrepancy between theoretically predicted and actually observed contractual forms.
    \item \textcite{Lazear2007} revisited this question: ``there must be some benefit of these contracts that outweighs their apparent costs.''
\end{itemize}


\section{Contract Theory \& Incentive Schemes}

\subsection*{The Puzzle of Contract Simplicity}

\textbf{\textcite{Prendergast1999}: \textit{The Provision of Incentives in Firms}}
\begin{itemize}
    \item \textbf{Documents the puzzle:} Real contracts are far simpler than theory predicts --- but \emph{does not resolve} it.
    \item \textbf{Contribution:} Highlights the disconnect between complex theoretical prescriptions and standard contractual practice.
    \item \textbf{Methodological limit:} Survey-based, no causal identification of \emph{why} simplicity prevails.
\end{itemize}

\medskip

\textbf{\textcite{Lazear2007}: \textit{Personnel Economics}}
\begin{itemize}
    \item \textbf{Sharpens the puzzle:} Argues there must be benefits to simple contracts that classical theory misses.
    \item \textbf{Contribution:} Bridges the gap between theoretical models and applied personnel economics.
    \item \textbf{Open question:} What \emph{exactly} are these benefits? --- This is where \textcite{Herweg2010} enters.
\end{itemize}


\section{Theoretical Framework}

\subsection*{A Behavioural Explanation}

\textcite{Herweg2010} address this puzzle by introducing \textbf{expectation-based loss aversion} into the standard principal--agent model with moral hazard, building on \textcite{Koszegi2006}.

\bigskip

\textbf{The K\H{o}szegi--Rabin Concept:}
\begin{itemize}
    \item People feel losses more strongly than gains of the same magnitude (Prospect Theory, \cite{Kahneman1979}).
    \item Expectation-based loss aversion means that the agent measures their utility not only by the actual outcome, but also by whether the outcome is better or worse than their expectations.
\end{itemize}

Why this matters for your story
Your current “Key Insight” block says the agent dislikes uncertainty because it increases the probability of falling short of the expected value. That is slightly misleading: in Kőszegi–Rabin, the agent does not just compare to the mean; he compares to the set of possible wages he expected. If the contract promises five different possible wages, the reference distribution has five mass points, and every realized wage is compared against the four other (higher) possible wages, generating loss utility. This is exactly why binary schemes are optimal in HMW: a binary contract ( wbase,wbonus ) creates a reference distribution with only two mass points, minimizing the number of “comparison points” against which the agent can feel disappointment, while still providing effort incentives.\\
\textbf{Stochastic loss aversion:} The agent compares his realized wage against \emph{every possible} wage in the reference distribution. Because the gain-loss function is kinked at zero, any wage dispersion in the reference lottery generates first-order psychological costs.

\begin{block}{Key Insight}
    The agent is risk-averse about stochastic reference points,\footnote{In equilibrium, the reference point coincides with rational expectations over wages.} because any additional uncertainty increases the probability that their payoff will fall short of the expected value ($\rightarrow$ loss aversion).
\end{block}


\section{Empirical Evidence}

\subsection*{Do Expectations Really Matter? Empirical Support}

The behavioural foundations of the model by \textcite{Herweg2010} rest on a well-established empirical foundation:

\medskip

\begin{itemize}
    \item \textbf{Field Evidence (Game Shows):} \textcite{Post2008} find that contestants' risk attitudes in \textit{Deal or No Deal} shift systematically with prior expectations.
    \item \textbf{Labour Supply (Cabdrivers):} \textcite{Crawford2011} shows that a model with expectation-based targets significantly outperforms the neoclassical benchmark in predicting cab drivers' effort.
    \item \textbf{Lab Evidence (Real-Effort):} \textcite{Abeler2011} demonstrates that directly manipulating subjects' expectations shifts effort provision in the direction predicted by loss aversion models.
\end{itemize}

\medskip

\small\alert{Note:} The central prediction, that binary schemes are optimal, has not yet been directly tested. This is addressed in the second presentation.


\section{Model Results \& Outlook: The Bridge}

\subsection*{Main Result: Why Binary Schemes?}

\textbf{\textcite{Herweg2010} identify a trade-off:}\newline
Two opposing effects of adding wage levels:
\begin{enumerate}
    \item \textbf{Incentive Gain:} More wage levels (complexity) provide better incentives. More wage levels $\Rightarrow$ finer performance signals $\Rightarrow$ higher marginal returns to effort.
    \item \textbf{Loss Aversion Cost:} More wage levels increase the probability of ``falling short'' of an expectation, creating ``psychological costs.'' More wage levels $\Rightarrow$ more reference points $\Rightarrow$ a higher probability of falling short of expectations.
\end{enumerate}

A binary scheme creates \textbf{exactly one} reference point --- minimising loss aversion costs while still providing meaningful incentives.

\begin{block}{The Optimality Theorem}
    When loss aversion is sufficiently high, the principal minimises these psychological costs by offering only two wages: a \textbf{base wage} and a \textbf{single bonus}.
\end{block}

\subsection*{From Theory to Experiment: Looking Ahead}

The theoretical results raise an important open question:

\bigskip

\begin{alertblock}{Research Gap}
    While theory establishes that binary schemes are optimal \emph{given} expectation-based loss aversion, it remains unclear how sensitive this result is to the \emph{way in which} expectations are formed.
\end{alertblock}

\textbf{Subsequent Theoretical Work:}\newline
\textcite{Herweg2015} extend the framework to dynamic settings, showing that expectation-based loss aversion can also explain \emph{inefficient renegotiation} of contracts --- further evidence that the mechanism has explanatory power beyond static moral hazard.

\medskip

\textbf{Proposed Experimental Extension} (building on \textcite{Abeler2011}):
\begin{itemize}
    \item Do subjects reveal a preference for binary schemes when their expectations are \emph{explicitly} set to a high level?
    \item Does the \emph{simplicity} of the contract itself function as a commitment device against expectation-based losses?
\end{itemize}

\bigskip

\small\textit{These questions will be addressed in the second presentation.}


\section*{References}
\addcontentsline{toc}{section}{References}
\printbibliography

\end{document}
